% Options for packages loaded elsewhere
\PassOptionsToPackage{unicode}{hyperref}
\PassOptionsToPackage{hyphens}{url}
\documentclass[
]{article}
\usepackage{xcolor}
\usepackage[margin=1in]{geometry}
\usepackage{amsmath,amssymb}
\setcounter{secnumdepth}{-\maxdimen} % remove section numbering
\usepackage{iftex}
\ifPDFTeX
  \usepackage[T1]{fontenc}
  \usepackage[utf8]{inputenc}
  \usepackage{textcomp} % provide euro and other symbols
\else % if luatex or xetex
  \usepackage{unicode-math} % this also loads fontspec
  \defaultfontfeatures{Scale=MatchLowercase}
  \defaultfontfeatures[\rmfamily]{Ligatures=TeX,Scale=1}
\fi
\usepackage{lmodern}
\ifPDFTeX\else
  % xetex/luatex font selection
\fi
% Use upquote if available, for straight quotes in verbatim environments
\IfFileExists{upquote.sty}{\usepackage{upquote}}{}
\IfFileExists{microtype.sty}{% use microtype if available
  \usepackage[]{microtype}
  \UseMicrotypeSet[protrusion]{basicmath} % disable protrusion for tt fonts
}{}
\makeatletter
\@ifundefined{KOMAClassName}{% if non-KOMA class
  \IfFileExists{parskip.sty}{%
    \usepackage{parskip}
  }{% else
    \setlength{\parindent}{0pt}
    \setlength{\parskip}{6pt plus 2pt minus 1pt}}
}{% if KOMA class
  \KOMAoptions{parskip=half}}
\makeatother
\usepackage{color}
\usepackage{fancyvrb}
\newcommand{\VerbBar}{|}
\newcommand{\VERB}{\Verb[commandchars=\\\{\}]}
\DefineVerbatimEnvironment{Highlighting}{Verbatim}{commandchars=\\\{\}}
% Add ',fontsize=\small' for more characters per line
\usepackage{framed}
\definecolor{shadecolor}{RGB}{248,248,248}
\newenvironment{Shaded}{\begin{snugshade}}{\end{snugshade}}
\newcommand{\AlertTok}[1]{\textcolor[rgb]{0.94,0.16,0.16}{#1}}
\newcommand{\AnnotationTok}[1]{\textcolor[rgb]{0.56,0.35,0.01}{\textbf{\textit{#1}}}}
\newcommand{\AttributeTok}[1]{\textcolor[rgb]{0.13,0.29,0.53}{#1}}
\newcommand{\BaseNTok}[1]{\textcolor[rgb]{0.00,0.00,0.81}{#1}}
\newcommand{\BuiltInTok}[1]{#1}
\newcommand{\CharTok}[1]{\textcolor[rgb]{0.31,0.60,0.02}{#1}}
\newcommand{\CommentTok}[1]{\textcolor[rgb]{0.56,0.35,0.01}{\textit{#1}}}
\newcommand{\CommentVarTok}[1]{\textcolor[rgb]{0.56,0.35,0.01}{\textbf{\textit{#1}}}}
\newcommand{\ConstantTok}[1]{\textcolor[rgb]{0.56,0.35,0.01}{#1}}
\newcommand{\ControlFlowTok}[1]{\textcolor[rgb]{0.13,0.29,0.53}{\textbf{#1}}}
\newcommand{\DataTypeTok}[1]{\textcolor[rgb]{0.13,0.29,0.53}{#1}}
\newcommand{\DecValTok}[1]{\textcolor[rgb]{0.00,0.00,0.81}{#1}}
\newcommand{\DocumentationTok}[1]{\textcolor[rgb]{0.56,0.35,0.01}{\textbf{\textit{#1}}}}
\newcommand{\ErrorTok}[1]{\textcolor[rgb]{0.64,0.00,0.00}{\textbf{#1}}}
\newcommand{\ExtensionTok}[1]{#1}
\newcommand{\FloatTok}[1]{\textcolor[rgb]{0.00,0.00,0.81}{#1}}
\newcommand{\FunctionTok}[1]{\textcolor[rgb]{0.13,0.29,0.53}{\textbf{#1}}}
\newcommand{\ImportTok}[1]{#1}
\newcommand{\InformationTok}[1]{\textcolor[rgb]{0.56,0.35,0.01}{\textbf{\textit{#1}}}}
\newcommand{\KeywordTok}[1]{\textcolor[rgb]{0.13,0.29,0.53}{\textbf{#1}}}
\newcommand{\NormalTok}[1]{#1}
\newcommand{\OperatorTok}[1]{\textcolor[rgb]{0.81,0.36,0.00}{\textbf{#1}}}
\newcommand{\OtherTok}[1]{\textcolor[rgb]{0.56,0.35,0.01}{#1}}
\newcommand{\PreprocessorTok}[1]{\textcolor[rgb]{0.56,0.35,0.01}{\textit{#1}}}
\newcommand{\RegionMarkerTok}[1]{#1}
\newcommand{\SpecialCharTok}[1]{\textcolor[rgb]{0.81,0.36,0.00}{\textbf{#1}}}
\newcommand{\SpecialStringTok}[1]{\textcolor[rgb]{0.31,0.60,0.02}{#1}}
\newcommand{\StringTok}[1]{\textcolor[rgb]{0.31,0.60,0.02}{#1}}
\newcommand{\VariableTok}[1]{\textcolor[rgb]{0.00,0.00,0.00}{#1}}
\newcommand{\VerbatimStringTok}[1]{\textcolor[rgb]{0.31,0.60,0.02}{#1}}
\newcommand{\WarningTok}[1]{\textcolor[rgb]{0.56,0.35,0.01}{\textbf{\textit{#1}}}}
\usepackage{graphicx}
\makeatletter
\newsavebox\pandoc@box
\newcommand*\pandocbounded[1]{% scales image to fit in text height/width
  \sbox\pandoc@box{#1}%
  \Gscale@div\@tempa{\textheight}{\dimexpr\ht\pandoc@box+\dp\pandoc@box\relax}%
  \Gscale@div\@tempb{\linewidth}{\wd\pandoc@box}%
  \ifdim\@tempb\p@<\@tempa\p@\let\@tempa\@tempb\fi% select the smaller of both
  \ifdim\@tempa\p@<\p@\scalebox{\@tempa}{\usebox\pandoc@box}%
  \else\usebox{\pandoc@box}%
  \fi%
}
% Set default figure placement to htbp
\def\fps@figure{htbp}
\makeatother
\setlength{\emergencystretch}{3em} % prevent overfull lines
\providecommand{\tightlist}{%
  \setlength{\itemsep}{0pt}\setlength{\parskip}{0pt}}
\usepackage{booktabs}
\usepackage{longtable}
\usepackage{array}
\usepackage{multirow}
\usepackage{wrapfig}
\usepackage{float}
\usepackage{colortbl}
\usepackage{pdflscape}
\usepackage{tabularx}
\usepackage{xltabular}
\usepackage{threeparttable}
\usepackage{threeparttablex}
\usepackage[normalem]{ulem}
\usepackage{makecell}
\usepackage{xcolor}
\usepackage{bookmark}
\IfFileExists{xurl.sty}{\usepackage{xurl}}{} % add URL line breaks if available
\urlstyle{same}
\hypersetup{
  pdftitle={Data Wrangling (Data Preprocessing)},
  pdfauthor={Tung Duc Nguyen, Hoang Nguyen},
  hidelinks,
  pdfcreator={LaTeX via pandoc}}

\title{Data Wrangling (Data Preprocessing)}
\usepackage{etoolbox}
\makeatletter
\providecommand{\subtitle}[1]{% add subtitle to \maketitle
  \apptocmd{\@title}{\par {\large #1 \par}}{}{}
}
\makeatother
\subtitle{Practical assessment 1}
\author{Tung Duc Nguyen, Hoang Nguyen}
\date{}

\begin{document}
\maketitle

\subsection{\texorpdfstring{\textbf{Setup}}{Setup}}\label{setup}

\begin{Shaded}
\begin{Highlighting}[]
\CommentTok{\# Loading necessary packages:}
\FunctionTok{library}\NormalTok{(kableExtra)}
\FunctionTok{library}\NormalTok{(magrittr)}
\FunctionTok{library}\NormalTok{(readxl)}
\FunctionTok{library}\NormalTok{(dplyr)}
\end{Highlighting}
\end{Shaded}

\begin{verbatim}
## 
## Attaching package: 'dplyr'
\end{verbatim}

\begin{verbatim}
## The following object is masked from 'package:kableExtra':
## 
##     group_rows
\end{verbatim}

\begin{verbatim}
## The following objects are masked from 'package:stats':
## 
##     filter, lag
\end{verbatim}

\begin{verbatim}
## The following objects are masked from 'package:base':
## 
##     intersect, setdiff, setequal, union
\end{verbatim}

\begin{Shaded}
\begin{Highlighting}[]
\CommentTok{\#citation("dplyr")}
\end{Highlighting}
\end{Shaded}

\subsection{\texorpdfstring{\textbf{Student names, numbers and
percentage of
contributions}}{Student names, numbers and percentage of contributions}}\label{student-names-numbers-and-percentage-of-contributions}

\begin{table}
\centering
\caption{\label{tab:unnamed-chunk-2}Group information}
\centering
\begin{tabular}[t]{l|l|l}
\hline
Student name & Student number & Percentage of contribution\\
\hline
Tung Duc Nguyen & 4275232 & 50\\
\hline
Hoang Nguyen & 4280335 & 50\\
\hline
\end{tabular}
\end{table}

\subsection{\texorpdfstring{\textbf{Data
Description}}{Data Description}}\label{data-description}

Data set: Bail Justice and bail applications by outcome, age group and
Indigenous status --- data ending March 2026

Source: Crime Statistics Agency (CSA) Victoria, published via the
Victorian Government's open data platform
(\url{https://discover.data.vic.gov.au/dataset/bail-statistics})

The data set records the outcomes of bail applications heard by a Bail
Justice in Victoria, organized by age group and Indigenous status,
reported as 12-month totals ending each March from 2017 to 2026.

Variables:

\begin{itemize}
\tightlist
\item
  Year -- the year the 12-month reporting period ends
\item
  Year ending -- the month the period ends
\item
  Data Source -- the agency/court the data was recorded by
\item
  Outcome -- the result of the bail application (Bail granted, Bail
  refused, No application)
\item
  Age Group -- Adult 18 years and over, Child under 18 years, or Total
\item
  Indigenous Status -- Aboriginal and/or Torres Strait Islander,
  Non-Indigenous, or Total People
\item
  Record Count -- the number of people/applications matching that
  combination of Outcome, Age Group and Indigenous Status for that year
\end{itemize}

\subsection{\texorpdfstring{\textbf{Read/Import
Data}}{Read/Import Data}}\label{readimport-data}

\begin{Shaded}
\begin{Highlighting}[]
\CommentTok{\# Import the data, provide your R codes here.}
\FunctionTok{library}\NormalTok{(readxl)}
\NormalTok{bail\_justice\_data }\OtherTok{\textless{}{-}} \FunctionTok{read\_xlsx}\NormalTok{(}\StringTok{"Bail\_Justice\_Data\_Tables\_Bail\_Visualisation\_Year\_Ending\_March\_2026.xlsx"}\NormalTok{, }\AttributeTok{sheet =} \StringTok{"Table 01"}\NormalTok{)}
\FunctionTok{class}\NormalTok{(bail\_justice\_data)}
\end{Highlighting}
\end{Shaded}

\begin{verbatim}
## [1] "tbl_df"     "tbl"        "data.frame"
\end{verbatim}

First we ran the library(readxl) to load the package ``readxl'' into our
session so we can use its function: ``read\_xlsx'' to read our dataset
from an excel file.

Then we ran ``read\_xlsx'' to read our file
``Bail\_Justice\_Data\_Tables\_Bail\_Visualisation\_Year\_Ending\_March\_2026.xlsx'',
specifically sheet ``Table 01'' as that is where the data table is,
while other sheets contain unusable data. We assigned the result to a
new object named: bail\_justice\_data, which is now a dataframe in our R
environment. To double check that fact, we also ran the class function
on it.

\subsection{\texorpdfstring{\textbf{Inspect and
Understand}}{Inspect and Understand}}\label{inspect-and-understand}

\begin{Shaded}
\begin{Highlighting}[]
\CommentTok{\# Check the overall structure and object type of the dataset}
\FunctionTok{str}\NormalTok{(bail\_justice\_data)}
\end{Highlighting}
\end{Shaded}

\begin{verbatim}
## tibble [159 x 7] (S3: tbl_df/tbl/data.frame)
##  $ Year             : num [1:159] 2026 2026 2026 2026 2026 ...
##  $ Year ending      : chr [1:159] "March" "March" "March" "March" ...
##  $ Data Source      : chr [1:159] "Bail Justice" "Bail Justice" "Bail Justice" "Bail Justice" ...
##  $ Outcome          : chr [1:159] "Bail granted" "Bail granted" "Bail granted" "Bail granted" ...
##  $ Age Group        : chr [1:159] "Adult 18 years and over" "Adult 18 years and over" "Adult 18 years and over" "Child under 18 years" ...
##  $ Indigenous Status: chr [1:159] "Aboriginal and/or Torres Strait Islander" "Non-Indigenous" "Total People" "Aboriginal and/or Torres Strait Islander" ...
##  $ Record Count     : num [1:159] 77 43 120 2 15 17 79 58 137 855 ...
\end{verbatim}

\begin{Shaded}
\begin{Highlighting}[]
\FunctionTok{class}\NormalTok{(bail\_justice\_data)}
\end{Highlighting}
\end{Shaded}

\begin{verbatim}
## [1] "tbl_df"     "tbl"        "data.frame"
\end{verbatim}

Firstly for the inspection, we ran str() to get an overview of the
dataset's structure, which shows us the total number of observations and
variables, and for each variables it also shows the data type with a
preview of the first few values. Additionally, we also ran class() to
confirm that the dataset is a dataframe.

\begin{Shaded}
\begin{Highlighting}[]
\CommentTok{\# Preview the first and last rows of the dataset}
\FunctionTok{head}\NormalTok{(bail\_justice\_data)}
\end{Highlighting}
\end{Shaded}

\begin{verbatim}
## # A tibble: 6 x 7
##    Year `Year ending` `Data Source` Outcome      `Age Group` `Indigenous Status`
##   <dbl> <chr>         <chr>         <chr>        <chr>       <chr>              
## 1  2026 March         Bail Justice  Bail granted Adult 18 y~ Aboriginal and/or ~
## 2  2026 March         Bail Justice  Bail granted Adult 18 y~ Non-Indigenous     
## 3  2026 March         Bail Justice  Bail granted Adult 18 y~ Total People       
## 4  2026 March         Bail Justice  Bail granted Child unde~ Aboriginal and/or ~
## 5  2026 March         Bail Justice  Bail granted Child unde~ Non-Indigenous     
## 6  2026 March         Bail Justice  Bail granted Child unde~ Total People       
## # i 1 more variable: `Record Count` <dbl>
\end{verbatim}

\begin{Shaded}
\begin{Highlighting}[]
\FunctionTok{head}\NormalTok{(bail\_justice\_data, }\AttributeTok{n =} \DecValTok{10}\NormalTok{)}
\end{Highlighting}
\end{Shaded}

\begin{verbatim}
## # A tibble: 10 x 7
##     Year `Year ending` `Data Source` Outcome     `Age Group` `Indigenous Status`
##    <dbl> <chr>         <chr>         <chr>       <chr>       <chr>              
##  1  2026 March         Bail Justice  Bail grant~ Adult 18 y~ Aboriginal and/or ~
##  2  2026 March         Bail Justice  Bail grant~ Adult 18 y~ Non-Indigenous     
##  3  2026 March         Bail Justice  Bail grant~ Adult 18 y~ Total People       
##  4  2026 March         Bail Justice  Bail grant~ Child unde~ Aboriginal and/or ~
##  5  2026 March         Bail Justice  Bail grant~ Child unde~ Non-Indigenous     
##  6  2026 March         Bail Justice  Bail grant~ Child unde~ Total People       
##  7  2026 March         Bail Justice  Bail grant~ Total       Aboriginal and/or ~
##  8  2026 March         Bail Justice  Bail grant~ Total       Non-Indigenous     
##  9  2026 March         Bail Justice  Bail grant~ Total       Total People       
## 10  2026 March         Bail Justice  Bail refus~ Adult 18 y~ Aboriginal and/or ~
## # i 1 more variable: `Record Count` <dbl>
\end{verbatim}

\begin{Shaded}
\begin{Highlighting}[]
\FunctionTok{tail}\NormalTok{(bail\_justice\_data)}
\end{Highlighting}
\end{Shaded}

\begin{verbatim}
## # A tibble: 6 x 7
##    Year `Year ending` `Data Source` Outcome      `Age Group` `Indigenous Status`
##   <dbl> <chr>         <chr>         <chr>        <chr>       <chr>              
## 1  2017 March         Bail Justice  Bail granted Adult 18 y~ Total People       
## 2  2017 March         Bail Justice  Bail granted Child unde~ Total People       
## 3  2017 March         Bail Justice  Bail granted Total       Total People       
## 4  2017 March         Bail Justice  Bail refused Adult 18 y~ Total People       
## 5  2017 March         Bail Justice  Bail refused Child unde~ Total People       
## 6  2017 March         Bail Justice  Bail refused Total       Total People       
## # i 1 more variable: `Record Count` <dbl>
\end{verbatim}

We ran head() to preview the first 6 rows of our dataset by default and
then ran it again with n = 10 to view the first 10 rows for a slightly
larger data sample. Afterwards we ran tail() to preview the last six
rows of the dataset to check that both the head and tail had been
imported correctly.

\begin{Shaded}
\begin{Highlighting}[]
\CommentTok{\# Check for missing values}
\FunctionTok{colSums}\NormalTok{(}\FunctionTok{is.na}\NormalTok{(bail\_justice\_data))}
\end{Highlighting}
\end{Shaded}

\begin{verbatim}
##              Year       Year ending       Data Source           Outcome 
##                 0                 0                 0                 0 
##         Age Group Indigenous Status      Record Count 
##                 0                 0                 0
\end{verbatim}

We used colSums(is.na()) to check for missing values in each individual
column. The table presented zero count of missing values so there are no
missing values in the data set.

To achieve this, We referred to a GeeksforGeeks tutorial on finding and
counting missing values in R to confirm the correct functions to use for
this check
(\url{https://www.geeksforgeeks.org/r-language/how-to-find-and-count-missing-values-in-r-dataframe/})

\begin{Shaded}
\begin{Highlighting}[]
\CommentTok{\# Check the data type of individual variables}
\FunctionTok{class}\NormalTok{(bail\_justice\_data}\SpecialCharTok{$}\NormalTok{Outcome)}
\end{Highlighting}
\end{Shaded}

\begin{verbatim}
## [1] "character"
\end{verbatim}

\begin{Shaded}
\begin{Highlighting}[]
\FunctionTok{class}\NormalTok{(bail\_justice\_data}\SpecialCharTok{$}\StringTok{\textasciigrave{}}\AttributeTok{Record Count}\StringTok{\textasciigrave{}}\NormalTok{)}
\end{Highlighting}
\end{Shaded}

\begin{verbatim}
## [1] "numeric"
\end{verbatim}

\begin{Shaded}
\begin{Highlighting}[]
\FunctionTok{typeof}\NormalTok{(bail\_justice\_data}\SpecialCharTok{$}\NormalTok{Year)}
\end{Highlighting}
\end{Shaded}

\begin{verbatim}
## [1] "double"
\end{verbatim}

\begin{Shaded}
\begin{Highlighting}[]
\FunctionTok{typeof}\NormalTok{(bail\_justice\_data}\SpecialCharTok{$}\StringTok{\textasciigrave{}}\AttributeTok{Record Count}\StringTok{\textasciigrave{}}\NormalTok{)}
\end{Highlighting}
\end{Shaded}

\begin{verbatim}
## [1] "double"
\end{verbatim}

\begin{Shaded}
\begin{Highlighting}[]
\FunctionTok{is.numeric}\NormalTok{(bail\_justice\_data}\SpecialCharTok{$}\NormalTok{Year)}
\end{Highlighting}
\end{Shaded}

\begin{verbatim}
## [1] TRUE
\end{verbatim}

\begin{Shaded}
\begin{Highlighting}[]
\FunctionTok{is.numeric}\NormalTok{(bail\_justice\_data}\SpecialCharTok{$}\StringTok{\textasciigrave{}}\AttributeTok{Record Count}\StringTok{\textasciigrave{}}\NormalTok{)}
\end{Highlighting}
\end{Shaded}

\begin{verbatim}
## [1] TRUE
\end{verbatim}

\begin{Shaded}
\begin{Highlighting}[]
\FunctionTok{is.numeric}\NormalTok{(bail\_justice\_data}\SpecialCharTok{$}\StringTok{\textasciigrave{}}\AttributeTok{Data Source}\StringTok{\textasciigrave{}}\NormalTok{)}
\end{Highlighting}
\end{Shaded}

\begin{verbatim}
## [1] FALSE
\end{verbatim}

We ran class() on the Outcome and Record Count columns specifically to
check the data type of each and then used typeof() on Year and Record
Count to check how R is storing these numeric variables. Finally, we
used is.numeric() to explicitly test whether Year, Record Count, and
Data Source are numeric variables. This confirmed that Year and Record
Count are numeric (returning TRUE), while Data Source is not (returning
FALSE), since it's a text variable rather than a numeric one.

\begin{Shaded}
\begin{Highlighting}[]
\CommentTok{\# Check the attributes (metadata) of the dataset}
\FunctionTok{attributes}\NormalTok{(bail\_justice\_data)}
\end{Highlighting}
\end{Shaded}

\begin{verbatim}
## $names
## [1] "Year"              "Year ending"       "Data Source"      
## [4] "Outcome"           "Age Group"         "Indigenous Status"
## [7] "Record Count"     
## 
## $row.names
##   [1]   1   2   3   4   5   6   7   8   9  10  11  12  13  14  15  16  17  18
##  [19]  19  20  21  22  23  24  25  26  27  28  29  30  31  32  33  34  35  36
##  [37]  37  38  39  40  41  42  43  44  45  46  47  48  49  50  51  52  53  54
##  [55]  55  56  57  58  59  60  61  62  63  64  65  66  67  68  69  70  71  72
##  [73]  73  74  75  76  77  78  79  80  81  82  83  84  85  86  87  88  89  90
##  [91]  91  92  93  94  95  96  97  98  99 100 101 102 103 104 105 106 107 108
## [109] 109 110 111 112 113 114 115 116 117 118 119 120 121 122 123 124 125 126
## [127] 127 128 129 130 131 132 133 134 135 136 137 138 139 140 141 142 143 144
## [145] 145 146 147 148 149 150 151 152 153 154 155 156 157 158 159
## 
## $class
## [1] "tbl_df"     "tbl"        "data.frame"
\end{verbatim}

We ran attributes() on our data frame to show the column names, the row
names, and the class of the object, which is useful for confirming the
overall shape and labeling of the dataset beyond just the individual
variable types.

\begin{Shaded}
\begin{Highlighting}[]
\CommentTok{\# Convert categorical variables from character to factor}
\NormalTok{bail\_justice\_data[}\FunctionTok{c}\NormalTok{(}\StringTok{"Outcome"}\NormalTok{, }\StringTok{"Age Group"}\NormalTok{, }\StringTok{"Indigenous Status"}\NormalTok{, }\StringTok{"Data Source"}\NormalTok{)] }\OtherTok{\textless{}{-}} \FunctionTok{lapply}\NormalTok{(bail\_justice\_data[}\FunctionTok{c}\NormalTok{(}\StringTok{"Outcome"}\NormalTok{, }\StringTok{"Age Group"}\NormalTok{, }\StringTok{"Indigenous Status"}\NormalTok{, }\StringTok{"Data Source"}\NormalTok{)], as.factor)}
\FunctionTok{str}\NormalTok{(bail\_justice\_data)}
\end{Highlighting}
\end{Shaded}

\begin{verbatim}
## tibble [159 x 7] (S3: tbl_df/tbl/data.frame)
##  $ Year             : num [1:159] 2026 2026 2026 2026 2026 ...
##  $ Year ending      : chr [1:159] "March" "March" "March" "March" ...
##  $ Data Source      : Factor w/ 1 level "Bail Justice": 1 1 1 1 1 1 1 1 1 1 ...
##  $ Outcome          : Factor w/ 3 levels "Bail granted",..: 1 1 1 1 1 1 1 1 1 2 ...
##  $ Age Group        : Factor w/ 3 levels "Adult 18 years and over",..: 1 1 1 2 2 2 3 3 3 1 ...
##  $ Indigenous Status: Factor w/ 3 levels "Aboriginal and/or Torres Strait Islander",..: 1 2 3 1 2 3 1 2 3 1 ...
##  $ Record Count     : num [1:159] 77 43 120 2 15 17 79 58 137 855 ...
\end{verbatim}

\begin{Shaded}
\begin{Highlighting}[]
\FunctionTok{levels}\NormalTok{(bail\_justice\_data}\SpecialCharTok{$}\StringTok{\textasciigrave{}}\AttributeTok{Age Group}\StringTok{\textasciigrave{}}\NormalTok{)}
\end{Highlighting}
\end{Shaded}

\begin{verbatim}
## [1] "Adult 18 years and over" "Child under 18 years"   
## [3] "Total"
\end{verbatim}

Since Outcome, Age Group, Indigenous Status, and Data Source are all
categorical variables that were imported as character type, we used
lapply() to apply the as.factor() function to all four columns and
reassign them back into the data frame. This converts each of them from
character to factor, which is more appropriate for categorical data. we
then ran str() again to confirm the four columns are now listed as
factors instead of characters, and ran levels() on the Age Group column
specifically to check what categories exist now that it has been
converted, confirming the conversion worked.

\begin{Shaded}
\begin{Highlighting}[]
\CommentTok{\# Convert Year from numeric to integer}
\NormalTok{bail\_justice\_data}\SpecialCharTok{$}\NormalTok{Year }\OtherTok{\textless{}{-}} \FunctionTok{as.integer}\NormalTok{(bail\_justice\_data}\SpecialCharTok{$}\NormalTok{Year)}
\NormalTok{bail\_justice\_data}\SpecialCharTok{$}\StringTok{\textasciigrave{}}\AttributeTok{Record Count}\StringTok{\textasciigrave{}} \OtherTok{\textless{}{-}} \FunctionTok{as.integer}\NormalTok{(bail\_justice\_data}\SpecialCharTok{$}\StringTok{\textasciigrave{}}\AttributeTok{Record Count}\StringTok{\textasciigrave{}}\NormalTok{)}
\FunctionTok{typeof}\NormalTok{(bail\_justice\_data}\SpecialCharTok{$}\NormalTok{Year)}
\end{Highlighting}
\end{Shaded}

\begin{verbatim}
## [1] "integer"
\end{verbatim}

\begin{Shaded}
\begin{Highlighting}[]
\FunctionTok{typeof}\NormalTok{(bail\_justice\_data}\SpecialCharTok{$}\StringTok{\textasciigrave{}}\AttributeTok{Record Count}\StringTok{\textasciigrave{}}\NormalTok{)}
\end{Highlighting}
\end{Shaded}

\begin{verbatim}
## [1] "integer"
\end{verbatim}

Finally, we converted the Year column and the Record Count column from
numeric (double) to integer using as.integer(), since they are really
whole-numbers rather than continuous measurements, therefore integer is
the more appropriate and accurate type for them we then ran typeof()
again to confirm the conversions were successful and it now shows as
``integer'' rather than ``double''.

\subsection{\texorpdfstring{\textbf{Summary Statistics
}}{Summary Statistics }}\label{summary-statistics}

\begin{Shaded}
\begin{Highlighting}[]
\CommentTok{\# Summary statistics, provide R codes here.}
\CommentTok{\#Filter out "Total" (Age Group) and "Total People" (Indigenous Status)}
\NormalTok{bail\_justice\_data }\OtherTok{\textless{}{-}}\NormalTok{ bail\_justice\_data }\SpecialCharTok{\%\textgreater{}\%}
\FunctionTok{filter}\NormalTok{(}\StringTok{\textasciigrave{}}\AttributeTok{Age Group}\StringTok{\textasciigrave{}} \SpecialCharTok{!=} \StringTok{"Total"}\NormalTok{, }\StringTok{\textasciigrave{}}\AttributeTok{Indigenous Status}\StringTok{\textasciigrave{}} \SpecialCharTok{!=} \StringTok{"Total People"}\NormalTok{)}
\NormalTok{summary\_stats }\OtherTok{\textless{}{-}}\NormalTok{ bail\_justice\_data }\SpecialCharTok{\%\textgreater{}\%}
  \FunctionTok{group\_by}\NormalTok{(Outcome) }\SpecialCharTok{\%\textgreater{}\%}
  \FunctionTok{summarise}\NormalTok{(}
    \AttributeTok{Mean =} \FunctionTok{mean}\NormalTok{(}\StringTok{\textasciigrave{}}\AttributeTok{Record Count}\StringTok{\textasciigrave{}}\NormalTok{, }\AttributeTok{na.rm =} \ConstantTok{TRUE}\NormalTok{),}
    \AttributeTok{Median =} \FunctionTok{median}\NormalTok{(}\StringTok{\textasciigrave{}}\AttributeTok{Record Count}\StringTok{\textasciigrave{}}\NormalTok{, }\AttributeTok{na.rm =} \ConstantTok{TRUE}\NormalTok{),}
    \AttributeTok{Q1 =} \FunctionTok{quantile}\NormalTok{(}\StringTok{\textasciigrave{}}\AttributeTok{Record Count}\StringTok{\textasciigrave{}}\NormalTok{, }\FloatTok{0.25}\NormalTok{, }\AttributeTok{na.rm =} \ConstantTok{TRUE}\NormalTok{),}
    \AttributeTok{Q3 =} \FunctionTok{quantile}\NormalTok{(}\StringTok{\textasciigrave{}}\AttributeTok{Record Count}\StringTok{\textasciigrave{}}\NormalTok{, }\FloatTok{0.75}\NormalTok{, }\AttributeTok{na.rm =} \ConstantTok{TRUE}\NormalTok{),}
    \AttributeTok{SD =} \FunctionTok{sd}\NormalTok{(}\StringTok{\textasciigrave{}}\AttributeTok{Record Count}\StringTok{\textasciigrave{}}\NormalTok{, }\AttributeTok{na.rm =} \ConstantTok{TRUE}\NormalTok{)}
\NormalTok{  )}
\NormalTok{summary\_stats}
\end{Highlighting}
\end{Shaded}

\begin{verbatim}
## # A tibble: 3 x 6
##   Outcome         Mean Median    Q1    Q3    SD
##   <fct>          <dbl>  <dbl> <dbl> <dbl> <dbl>
## 1 Bail granted    63.0   56    33.5  100.  44.7
## 2 Bail refused   480.   453   131.   788. 385. 
## 3 No application  73.1   50.5  24.8  120.  63.9
\end{verbatim}

Before generating summary statistics, we filtered out the `Total' (Age
Group) and `Total People' (Indigenous Status) rows, as they are
subtotals of the other categories rather than independent values, and
including them would double-count people and mix these less detailed
value into the calculation and reduces the quality of our data.

To start,we grouped the bail justice data by the categorical variable
``Outcome'' and calculated the summary statistics for the numeric
variable ``Record Count''. Then, We used mean() to find the average
record count, median() to find the middle value, quantile() to calculate
the first and third quartiles, and sd() to find the standard deviation
for each outcome. Additionally, We also used na.rm = TRUE so any missing
values would be ignored during the calculations.

From the output, we found that `Bail refused' has the highest mean
Record Count of approximately 479.92 and also the highest median of 453,
while having the highest standard deviation of approximately 384.78, all
of which shows that the Record Count for Bail refused varies
considerably more than the other outcomes. Additionally, `Bail granted'
has the lowest mean Record Count of approximately 62.96 and the lowest
standard deviation of approximately 44.66.''

\subsection{\texorpdfstring{\textbf{Subsetting}}{Subsetting}}\label{subsetting}

\begin{Shaded}
\begin{Highlighting}[]
\CommentTok{\# Subset your data and convert it to a matrix, provide R codes here.}
\NormalTok{bail\_subset }\OtherTok{\textless{}{-}}\NormalTok{ bail\_justice\_data[}\DecValTok{1}\SpecialCharTok{:}\DecValTok{10}\NormalTok{, ]}
\NormalTok{bail\_subset}
\end{Highlighting}
\end{Shaded}

\begin{verbatim}
## # A tibble: 10 x 7
##     Year `Year ending` `Data Source` Outcome     `Age Group` `Indigenous Status`
##    <int> <chr>         <fct>         <fct>       <fct>       <fct>              
##  1  2026 March         Bail Justice  Bail grant~ Adult 18 y~ Aboriginal and/or ~
##  2  2026 March         Bail Justice  Bail grant~ Adult 18 y~ Non-Indigenous     
##  3  2026 March         Bail Justice  Bail grant~ Child unde~ Aboriginal and/or ~
##  4  2026 March         Bail Justice  Bail grant~ Child unde~ Non-Indigenous     
##  5  2026 March         Bail Justice  Bail refus~ Adult 18 y~ Aboriginal and/or ~
##  6  2026 March         Bail Justice  Bail refus~ Adult 18 y~ Non-Indigenous     
##  7  2026 March         Bail Justice  Bail refus~ Child unde~ Aboriginal and/or ~
##  8  2026 March         Bail Justice  Bail refus~ Child unde~ Non-Indigenous     
##  9  2026 March         Bail Justice  No applica~ Adult 18 y~ Aboriginal and/or ~
## 10  2026 March         Bail Justice  No applica~ Adult 18 y~ Non-Indigenous     
## # i 1 more variable: `Record Count` <int>
\end{verbatim}

\begin{Shaded}
\begin{Highlighting}[]
\NormalTok{bail\_matrix }\OtherTok{\textless{}{-}} \FunctionTok{as.matrix}\NormalTok{(bail\_subset)}
\NormalTok{bail\_matrix}
\end{Highlighting}
\end{Shaded}

\begin{verbatim}
##       Year   Year ending Data Source    Outcome         
##  [1,] "2026" "March"     "Bail Justice" "Bail granted"  
##  [2,] "2026" "March"     "Bail Justice" "Bail granted"  
##  [3,] "2026" "March"     "Bail Justice" "Bail granted"  
##  [4,] "2026" "March"     "Bail Justice" "Bail granted"  
##  [5,] "2026" "March"     "Bail Justice" "Bail refused"  
##  [6,] "2026" "March"     "Bail Justice" "Bail refused"  
##  [7,] "2026" "March"     "Bail Justice" "Bail refused"  
##  [8,] "2026" "March"     "Bail Justice" "Bail refused"  
##  [9,] "2026" "March"     "Bail Justice" "No application"
## [10,] "2026" "March"     "Bail Justice" "No application"
##       Age Group                 Indigenous Status                         
##  [1,] "Adult 18 years and over" "Aboriginal and/or Torres Strait Islander"
##  [2,] "Adult 18 years and over" "Non-Indigenous"                          
##  [3,] "Child under 18 years"    "Aboriginal and/or Torres Strait Islander"
##  [4,] "Child under 18 years"    "Non-Indigenous"                          
##  [5,] "Adult 18 years and over" "Aboriginal and/or Torres Strait Islander"
##  [6,] "Adult 18 years and over" "Non-Indigenous"                          
##  [7,] "Child under 18 years"    "Aboriginal and/or Torres Strait Islander"
##  [8,] "Child under 18 years"    "Non-Indigenous"                          
##  [9,] "Adult 18 years and over" "Aboriginal and/or Torres Strait Islander"
## [10,] "Adult 18 years and over" "Non-Indigenous"                          
##       Record Count
##  [1,] " 77"       
##  [2,] " 43"       
##  [3,] "  2"       
##  [4,] " 15"       
##  [5,] "855"       
##  [6,] "548"       
##  [7,] " 41"       
##  [8,] "141"       
##  [9,] "185"       
## [10,] "112"
\end{verbatim}

\begin{Shaded}
\begin{Highlighting}[]
\FunctionTok{str}\NormalTok{(bail\_matrix)}
\end{Highlighting}
\end{Shaded}

\begin{verbatim}
##  chr [1:10, 1:7] "2026" "2026" "2026" "2026" "2026" "2026" "2026" "2026" ...
##  - attr(*, "dimnames")=List of 2
##   ..$ : NULL
##   ..$ : chr [1:7] "Year" "Year ending" "Data Source" "Outcome" ...
\end{verbatim}

\begin{Shaded}
\begin{Highlighting}[]
\FunctionTok{class}\NormalTok{(bail\_matrix)}
\end{Highlighting}
\end{Shaded}

\begin{verbatim}
## [1] "matrix" "array"
\end{verbatim}

\begin{Shaded}
\begin{Highlighting}[]
\FunctionTok{typeof}\NormalTok{(bail\_matrix)}
\end{Highlighting}
\end{Shaded}

\begin{verbatim}
## [1] "character"
\end{verbatim}

Firstly, we used {[}1:10, {]} to subset the first 10 observations from
the bail justice data while keeping all variables, then we used
as.matrix() to convert the subsetted data frame into a matrix.

From the output, we found that the matrix was a character matrix. This
happened because a matrix can only contain one data type, while a data
frame can contain different data types in different columns. Therefore,
since the subset contains categorical variables such as ``Outcome'',
``Age Group'' and ``Indigenous Status'' together with numeric variables
such as ``Year'' and ``Record Count'', R converts all values to
character so that the matrix can store them using one common data type.

\subsection{\texorpdfstring{\textbf{Create a new Data
Frame}}{Create a new Data Frame}}\label{create-a-new-data-frame}

\begin{Shaded}
\begin{Highlighting}[]
\CommentTok{\# Create a new data frame, provide R codes here.}

\NormalTok{student\_id }\OtherTok{\textless{}{-}} \DecValTok{1}\SpecialCharTok{:}\DecValTok{10}
\NormalTok{satisfaction }\OtherTok{\textless{}{-}} \FunctionTok{c}\NormalTok{(}\StringTok{"Low"}\NormalTok{, }\StringTok{"Medium"}\NormalTok{, }\StringTok{"High"}\NormalTok{, }\StringTok{"Medium"}\NormalTok{, }\StringTok{"Low"}\NormalTok{,}\StringTok{"High"}\NormalTok{, }\StringTok{"Medium"}\NormalTok{, }\StringTok{"High"}\NormalTok{, }\StringTok{"Low"}\NormalTok{, }\StringTok{"Medium"}\NormalTok{)}
\NormalTok{satisfaction }\OtherTok{\textless{}{-}} \FunctionTok{factor}\NormalTok{(satisfaction,}
                       \AttributeTok{levels =} \FunctionTok{c}\NormalTok{(}\StringTok{"Low"}\NormalTok{, }\StringTok{"Medium"}\NormalTok{, }\StringTok{"High"}\NormalTok{),}
                       \AttributeTok{ordered =} \ConstantTok{TRUE}\NormalTok{)}
\NormalTok{new\_data }\OtherTok{\textless{}{-}} \FunctionTok{data.frame}\NormalTok{(student\_id, satisfaction)}
\NormalTok{new\_data}
\end{Highlighting}
\end{Shaded}

\begin{verbatim}
##    student_id satisfaction
## 1           1          Low
## 2           2       Medium
## 3           3         High
## 4           4       Medium
## 5           5          Low
## 6           6         High
## 7           7       Medium
## 8           8         High
## 9           9          Low
## 10         10       Medium
\end{verbatim}

First, we created two vectors called ``student\_id'' and
``satisfaction'' with 10 observations: - The student\_id variable
contains integer values from 1 to 10. - The satisfaction variable
contains three ordinal categories: ``Low'', ``Medium'' and ``High''. We
used factor() with ordered = TRUE to convert satisfaction into an
ordered factor and set the levels from Low to Medium to High and then
used data.frame() to combine the two variables into a new data frame.

\begin{Shaded}
\begin{Highlighting}[]
\CommentTok{\# Check the structure and ordinal levels}

\FunctionTok{str}\NormalTok{(new\_data)}
\end{Highlighting}
\end{Shaded}

\begin{verbatim}
## 'data.frame':    10 obs. of  2 variables:
##  $ student_id  : int  1 2 3 4 5 6 7 8 9 10
##  $ satisfaction: Ord.factor w/ 3 levels "Low"<"Medium"<..: 1 2 3 2 1 3 2 3 1 2
\end{verbatim}

\begin{Shaded}
\begin{Highlighting}[]
\FunctionTok{levels}\NormalTok{(new\_data}\SpecialCharTok{$}\NormalTok{satisfaction)}
\end{Highlighting}
\end{Shaded}

\begin{verbatim}
## [1] "Low"    "Medium" "High"
\end{verbatim}

Next we used str() to check the structure of the new data frame and
levels() to check that the satisfaction variable was an ordered factor
with the levels in the correct order.

\begin{Shaded}
\begin{Highlighting}[]
\CommentTok{\# Create and add a numeric variable}
\NormalTok{hours\_studied }\OtherTok{\textless{}{-}} \FunctionTok{c}\NormalTok{(}\DecValTok{8}\NormalTok{, }\DecValTok{12}\NormalTok{, }\DecValTok{15}\NormalTok{, }\DecValTok{10}\NormalTok{, }\DecValTok{6}\NormalTok{, }\DecValTok{18}\NormalTok{, }\DecValTok{11}\NormalTok{, }\DecValTok{14}\NormalTok{, }\DecValTok{7}\NormalTok{, }\DecValTok{9}\NormalTok{)}
\NormalTok{new\_data }\OtherTok{\textless{}{-}} \FunctionTok{cbind}\NormalTok{(new\_data, hours\_studied)}
\NormalTok{new\_data}
\end{Highlighting}
\end{Shaded}

\begin{verbatim}
##    student_id satisfaction hours_studied
## 1           1          Low             8
## 2           2       Medium            12
## 3           3         High            15
## 4           4       Medium            10
## 5           5          Low             6
## 6           6         High            18
## 7           7       Medium            11
## 8           8         High            14
## 9           9          Low             7
## 10         10       Medium             9
\end{verbatim}

\begin{Shaded}
\begin{Highlighting}[]
\FunctionTok{str}\NormalTok{(new\_data)}
\end{Highlighting}
\end{Shaded}

\begin{verbatim}
## 'data.frame':    10 obs. of  3 variables:
##  $ student_id   : int  1 2 3 4 5 6 7 8 9 10
##  $ satisfaction : Ord.factor w/ 3 levels "Low"<"Medium"<..: 1 2 3 2 1 3 2 3 1 2
##  $ hours_studied: num  8 12 15 10 6 18 11 14 7 9
\end{verbatim}

Finally, we created another numeric vector called ``hours\_studied''
with 10 observations and used cbind() to add it to the existing data
frame, which outputed the final data frame containg 10 observations and
3 variables.

\subsection{\texorpdfstring{\textbf{References}}{References}}\label{references}

Bache SM and Wickham H (2026) \emph{magrittr: A Forward-Pipe Operator
for R}, version 2.0.5, R package,
\url{doi:10.32614/CRAN.package.magrittr}.

Wickham H and Bryan J (2026) \emph{readxl: Read Excel Files}, version
1.5.0, R package, \url{doi:10.32614/CRAN.package.readxl}.

Wickham H, François R, Henry L, Müller K and Vaughan D (2026)
\emph{dplyr: A Grammar of Data Manipulation}, version 1.2.1, R package,
\url{doi:10.32614/CRAN.package.dplyr}.

Zhu H (2026) \emph{kableExtra: Construct Complex Table with `kable' and
Pipe Syntax}, version 1.4.1, R package,
\url{doi:10.32614/CRAN.package.kableExtra}.

GeeksforGeeks (2022). How to find and count missing values in R
DataFrame. {[}online{]} GeeksforGeeks. Available at:
\url{https://www.geeksforgeeks.org/r-language/how-to-find-and-count-missing-values-in-r-dataframe/}
{[}Accessed 17 Aug.~2026{]}.

Crime Statistics Agency Data Tables - Bail statistics - DataVic. (2026).
Vic.Gov.Au.
\url{https://discover.data.vic.gov.au/dataset/bail-statistics}

TUNG1205 (2026). Commits ·
TUNG1205/bail-outcomes-in-victoria--age-and-indigenous-status-trends.
{[}online{]} GitHub. Available at:
\url{https://github.com/TUNG1205/Bail-Outcomes-in-Victoria--Age-and-Indigenous-Status-Trends/commits/main/}
{[}Accessed 20 Aug.~2026{]}.

\subsection{\texorpdfstring{\textbf{Presentation
}}{Presentation }}\label{presentation}

\href{https://rmit-arc.instructuremedia.com/embed/d76b8ee7-93f6-4c7a-8230-d336d4d8617c}{Tung's
Presentation}

\href{https://rmit-arc.instructuremedia.com/embed/a7c0b929-3974-4bcb-80c9-3f5f961197b7}{Hoang's
Presentation}

\end{document}
